\sect{वैशाख-शुक्ल-मोहिनी-एकादशी-माहात्म्यम्}
\label{sec:padma-vaishakha-shukla-mohini}

\ganapatyadiDhyanam
\padmaPuranam

\dnsub{अथ नामैकपञ्चाशत्तमोऽध्यायः॥५१}

\uvacha{युधिष्ठिर उवाच}

\twolineshloka
{वैशाख शुक्लपक्षे तु किं नामैकादशी भवेत्}
{किं फलं को विधिस्तत्र कथयस्व जनार्दन}% ॥१॥

\uvacha{श्रीकृष्ण उवाच}

\twolineshloka
{इदमेव पुरा पृष्टं रामचन्द्रेण धीमता}
{वसिष्ठं प्रति राजेन्द्र यत् त्वं मामनुपृच्छसि}% ॥२॥

\uvacha{राम उवाच}

\twolineshloka
{भगवञ्छ्रोतुमिच्छामि व्रतानामुत्तमं व्रतम्}
{सर्वपापक्षयकरं सर्वदुःखनिकृन्तनम्}% ॥३॥

\twolineshloka
{मया दुःखानि भुक्तानि सीताविरहजानि तु}
{ततोऽहं भयभीतोऽस्मि पृच्छामि त्वां महामुने}% ॥४॥

\uvacha{वसिष्ठ उवाच}

\twolineshloka
{साधु पृष्टं त्वया राम तवैषा नैष्ठिकी मतिः}
{त्वन्नामग्रहणेनैव पूतो भवति मानवः}% ॥५॥

\twolineshloka
{तथाऽपि कथयिष्यामि लोकानां हितकाम्यया}
{पवित्रं पावनानां च व्रतानामुत्तमं व्रतम्}% ॥६॥

\twolineshloka
{वैशाखस्य सिते पक्षे राम चैकादशी भवेत्}
{मोहिनी नाम सा प्रोक्ता सर्वपापहरापराः}% ॥७॥

\twolineshloka
{मोहजालात् प्रमुच्यन्ते पातकानां समूहतः}
{अस्या व्रत प्रभावेन सत्यं सत्यं वदाम्यहम्}% ॥८॥

\twolineshloka
{अतः कारणतो राम कर्तव्यैषा भवादृशैः}
{पातकानां क्षयकरी महादुःखविनाशिनी}% ॥९॥

\twolineshloka
{शृणुष्वैकमना राम कथां पापहरां पराम्}
{यस्याः श्रवणमात्रेण महापापं प्रणश्यति}% ॥१०॥

\twolineshloka
{सरस्वत्यास्तटे रम्ये पुरी भद्रावती शुभा}
{द्युतिमान्नाम नृपतिस्तत्र राज्यं करोति वै}% ॥११॥

\twolineshloka
{चन्द्रवंशोद्भवो नाम धृतिमान्सत्यसङ्गरः}
{तत्र वैश्यो निवसति धनधान्यसमृद्धिमान्}% ॥१२॥

\twolineshloka
{धनपाल इति ख्यातः पुण्यकर्मप्रवर्तकः}
{प्रपा कूप मठाराम तडाग गृहकारकः}% ॥१३॥

\twolineshloka
{विष्णुभक्तिरतः शान्तस्तस्यासन् पञ्चपुत्रकाः}
{सुमना द्युतिमांश्चैव मेधावी सुकृतस्तथा}% ॥१४॥

\twolineshloka
{पञ्चमो धृष्टबुद्धिश्च महापापरतः सदा}
{परस्त्रीसङ्गनिरतो विटगोष्ठी विशारदः}% ॥१५॥

\twolineshloka
{द्यूतादि व्यसनासक्तः परस्त्री रतिलालसः}
{न च देवार्चने बुद्धिर्नपितॄन्न द्विजान् प्रति}% ॥१६॥

\twolineshloka
{अन्यायवर्ती दुष्टात्मा पितुर्द्रव्यक्षयङ्करः}
{अभक्ष्यभक्षकः पापी सुरापाने रतः सदा}% ॥१७॥

\twolineshloka
{वेश्याकण्ठे क्षिप्तबाहुर्भ्रमन्दुष्टश्चतुष्पथे}
{पित्रा निष्कासितो गेहात् परित्यक्तश्च बान्धवैः}% ॥१८॥

\twolineshloka
{स्वदेहभूषणान्येव क्षयं नीतानि तेन वै}
{गणिकाभिः परित्यक्तो निन्दितश्च धनक्षयात्}% ॥१९॥

\twolineshloka
{ततश्चिन्तापरो जातो वस्त्रहीनः क्षुधार्दितः}
{किं करोमि क्वगच्छामि केनोपायेन जीव्यते}% ॥२०॥

\twolineshloka
{तस्करत्वं समारब्धं तत्रैव नगरे पितुः}
{गृहीतो राजपुरुषैर्मुक्तश्च पितृगौरवात्}% ॥२१॥

\twolineshloka
{पुनर्बद्धः पुनस्त्यक्तः पुनर्बद्धः ससम्भ्रमैः}
{धृष्टबुद्धिर्दुराचारो निबध्य निगडै र्दृढैः}% ॥२२॥

\twolineshloka
{कशाघातैस्ताडितश्च पीडितश्च पुनः पुनः}
{न स्थातव्यं हि मन्दात्मंस्त्वया मद्देशगोचरे}% ॥२३॥

\twolineshloka
{एवमुक्त्वा ततो राज्ञा मोचितो दृढबन्धनात्}
{निर्जगाम भयात् तस्य गतोऽसौ गहनं वनम्}% ॥२४॥

\twolineshloka
{क्षुत्तृषापीडितश्चायमितश्चेतश्च धावति}
{सिंहवन्निजघानासौ मृग शूकर चित्रलान्}% ॥२५॥

\twolineshloka
{आमिषाहार निरतो वने तिष्ठति सर्वदा}
{करे शरासनं कृत्वा निषङ्गं पृष्ठ सङ्गतम्}% ॥२६॥

\twolineshloka
{अरण्यचारिणो हन्ति पक्षिणश्च पदाचरन्}
{चकोरांश्च मयूरांश्च कङ्क तित्तिर मूषिकान्}% ॥२७॥

\twolineshloka
{एतानन्यान्हिनस्त्यन्धो धृष्टबुद्धिस्तु निर्घृणः}
{पूर्वजन्मकृतैः पापैर्निमग्नः पापकर्दमे}% ॥२८॥

\twolineshloka
{दुःखशोकसमाविष्टः पीड्यमानो दिवानिशम्}
{कौण्डिन्यस्याश्रमपदं प्राप्तः पुण्यागमात्क्वचित्}% ॥२९॥

\twolineshloka
{माधवे मासि जाह्नव्याः कृतस्नानं तपोधनम्}
{आससाद धृष्टबुद्धिः शोकभारेण पीडितः}% ॥३०॥

\twolineshloka
{तद्वस्त्रबिन्दुस्पर्शेन गतपापो हताशुभः}
{कौण्डिन्यस्याग्रतः स्थित्वा प्रत्युवाच कृताञ्जलि}% ॥३१॥

\uvacha{धृष्टबुद्धिरुवाच}

\twolineshloka
{भो भो ब्रह्मन् द्विजश्रेष्ठ दयां कृत्वा ममोपरि}
{येन पुण्यप्रभावेन मुक्तिर्भवति तद् वद}% ॥३२॥

\uvacha{कौण्डिन्य उवाच}

\twolineshloka
{शृणुष्वैकमनाभूत्वा येन पापक्षयस्तव}
{वैशाखस्य सिते पक्षे मोहिनी नाम विश्रुता}% ॥३३॥

\twolineshloka
{एकादशीव्रतं तस्याः कुरु मद्वाक्यनोदितः}
{मेरुतुल्यानि पापानि क्षयं गच्छन्ति देहिनाम्}% ॥३४॥

\twolineshloka
{बहुजन्मार्जितान्येषा मोहिनी समुपोषिता}
{इति वाक्यं मुनेः श्रुत्वा धृष्टबुद्धिः प्रसन्नधीः}% ॥३५॥

\twolineshloka
{व्रतं चकार विधिवत्कौण्डिन्यस्योपदेशतः}
{कृते व्रते नृपश्रेष्ठ गतपापो बभूव सः}% ॥३६॥

\twolineshloka
{दिव्यदेहस्ततो भूत्वा गरुडोपरिसंस्थितः}
{जगाम वैष्णवं लोकं सर्वोपद्रव वर्जितम्}% ॥३७॥

\twolineshloka
{इतीदृशं रामचन्द्र उत्तमं मोहिनी व्रतम्}
{नातः परतरं किञ्चित् त्रैलोक्ये सचराचरे}% ॥३८॥

\twolineshloka
{यज्ञादितीर्थदानानि कलां नार्हन्ति षोडशीम्}
{पठनाच्छ्रवणाद् राजन् गोसहस्रफलं लभेत्}% ॥३९॥

॥इति श्रीपाद्मे महापुराणे पञ्चपञ्चाशत्साहस्र्यां संहितायामुत्तरखण्डे उमापतिनारदसंवादान्तर्गतकृष्णयुधिष्ठिरसंवादे वैशाख-शुक्ल-मोहिनी-एकादशी-माहात्म्यं नाम नामैकपञ्चाशत्तमोऽध्यायः॥५१॥

\genMangalaShloka

\hyperref[sec:ekadashi_mahatmyam_padma_puranam]{\closesub}
\clearpage

